% v24: paper II への挿入用。v23 の固定族 Z_2 と復元写像の記号を用いる。
\section{自己反対点配置の固定点数と行列式曲線}\label{sec:self-opposite-21-v24}

\subsection{十一個の固定点と二次曲線上の五対}

最大の自己反対部分族について、点配置が有限被約になることと、
その二十一点に強制される幾何を同時に決定する。
ここでいう点配置上の対合は、反対合が左右の平面を同一視した後の
thin module の標準対応である。

\begin{theorem}[自己反対族の $11+10$ 分解]\label{thm:self-opposite-21-v24}
前節の $Z_2$ のある非空開部分では、左点スキームは相異なる二十一点からなる。
指定された反対合の誘導する点配置上の対合は十一点を固定し、
残る十点を五対に交換する。
この十点をちょうど含む非退化二次曲線が存在し、十一個の固定点はその外にある。

従って、この族の二十一点の monodromy は、適当なラベル付けの下で
\[
 \mathfrak S_{11}\times
 \bigl((\mathbb Z/2\mathbb Z)^5\rtimes\mathfrak S_5\bigr)
 \ \subset\ \mathfrak S_{21}
\]
に含まれる。この節では包含を示し、第\ref{sec:master-monodromy-v25}節で等号を証明する。
この非空開部分の同型類の次元は十一である。
\end{theorem}

まず線形代数から二次曲線を取り出す。
$D=\operatorname{diag}(1,-1,1)$、
$J=\operatorname{diag}(-1,-1,1,-1,-1,-1)$ と置く。
元の右枠では $R_j=\varepsilon_j J L_j$、
$(\varepsilon_1,\varepsilon_2,\varepsilon_3)=(1,-1,1)$ である。
右座標を $y=Dz$ と変更して左右を同一視すれば、以後は
\[
 R_i=J L_i
\]
と書ける。この座標変更は指定された反対合を変えるのではなく、
左右の平面の同一視を記述するものである。
基点での一致は
\[
 L_-(Dx)=(\operatorname{Sym}^2D)L_-(x)T,
 \qquad T=\operatorname{diag}(1,1,-1,-1,1,1,-1)
\]
からも直接確認できる。

\begin{lemma}[対称な復元 pairing]\label{lem:symmetric-pairing-v24}
復元写像の階数が四十八で、その核が可逆行列 $A$ で生成されるとする。
上の座標では $A$ は対称である。
$J$ の正固有空間に対応する行を $a$、負固有空間に対応する五行を $B$ と書くと、
\begin{equation}\label{eq:self-opposite-orthogonality-v24}
 a(x)A B(x)^t=0.
\end{equation}
\end{lemma}
\begin{proof}
復元方程式は
\[
 L_i A^t L_j^tJ=J L_j A L_i^t
\]
である。$i,j$ を交換し両側に $J$ を掛ければ、$A^t$ も解になる。
核が一次元なので $A^t=cA$、従って $c=\pm1$ である。
奇数次の可逆交代行列は存在しないので $c=1$ である。
正・負固有空間間の成分を比較すれば
\[
 a_i A B_j^t+a_j A B_i^t=0
\]
を得る。$i=j$ も含めて和を取れば所望の恒等式になる。
\end{proof}

\begin{proposition}[有限被約の場合の分解]\label{prop:11-plus-10-v24}
補題の仮定に加え、$a$ の三つの係数行が独立で、
左点スキーム $Z$ が有限被約であるとする。
このとき指定反対合の固定点はちょうど十一個である。
残る十点は
\[
 Q=V(q),\qquad q(x)=a(x)Aa(x)^t
\]
上にあり、固定点は $Q$ の外にある。
特に $q$ が非退化なら定理の配置を得る。
\end{proposition}
\begin{proof}
有限被約な maximal-minor スキームの各点で $L$ の階数は五である。
実際、階数四以下なら全六次小行列式の一次微分が消え、
零次元被約点を定めることができない。
従って各 $x\in Z$ に対して $L(x)^t u=0$ を満たす $[u]\in\mathbb P^5$ は一意である。

対称積の等式により、$z\mapsto (J L(z)A)^t u$ の像は
二次元空間 $\ker L(x)$ に入る。
この写像の階数が一以下なら、その核の射影化は直線を含み、
右点スキームも直線を含む。しかし右点スキームは $Z$ と同じなので、これは有限性に反する。
従って階数は二であり、thin module の左右射影は同型となる。
点配置上の対合 $\gamma$ は
\[
 L(\gamma(x))^t Ju=0
\]
で特徴付けられる。
$\gamma(x)=x$ であることは $Ju\sim u$ と同値である。
$a$ に基点がないので、$u$ が一次元の正固有空間に入ることはない。
よって固定点スキームは
\[
 T=\{x:\operatorname{rank}B(x)\leq4\}
\]
である。

恒等式 \eqref{eq:self-opposite-orthogonality-v24} により、$B$ の転置は
\[
 \mathcal O^5\longrightarrow
 E:=\ker\bigl(\mathcal O^7\xrightarrow{\,aA\,}\mathcal O(1)\bigr)(1)
\]
を定める。$a$ の係数行が独立だから
\[
 E\simeq\Omega^1_{\mathbb P^2}(2)\oplus\mathcal O(1)^4,
 \qquad c_1(E)=5H,\qquad c_2(E)=11H^2.
\]
$I_6(L)\subset I_5(B)$ より、$T$ は $Z$ の閉部分スキームなので有限被約である。
この五本の切断の退化 locus の局所 Hilbert--Burch 表示から
\[
 0\longrightarrow\mathcal O^5\longrightarrow E
   \longrightarrow\mathcal I_T(5)\longrightarrow0
\]
を得る。従って $\operatorname{length}T=c_2(E)=11$ である。
これは期待余次元二の場合のThom--Porteous公式とも一致する
\cite[\S14.4]{FultonIntersection}。

$x\in Z\setminus T$ では $B(x)$ の階数は五であり、
$L(x)$ の階数が五であることから $a(x)$ は $B(x)$ の行空間に入る。
式 \eqref{eq:self-opposite-orthogonality-v24} は直ちに $q(x)=0$ を与える。

最後に $x\in T\cap Q$ と仮定する。
$u$ は $J$ の負固有空間にあり、$B(x)^tu=0$ と書ける。
$k=Aa(x)^t$ は零でなく、$q(x)=0$ と直交恒等式から $k\in\ker L(x)$ である。
その恒等式を $\mathbb P^2$ の接方向 $v$ で微分し $u$ を掛けると
\[
 u^t(dL)_x(v)k=0
\]
を得る。
従って階数五の determinantal locus の法線微分
\[
 T_x\mathbb P^2\longrightarrow
 \operatorname{Hom}(\ker L(x),\operatorname{coker}L(x))
\]
は非零ベクトル $k$ を常に消し、階数は一以下となる。
標的は二次元なので、$Z$ は $x$ で被約零次元になれない。これは矛盾である。
ゆえに $T\cap Q=\varnothing$ であり、残る十点だけが $Q$ 上にある。
\end{proof}

\paragraph{同じ固定成分上での非空性の証明.}
有限被約性が $Z_2$ 上で実現することは、以下の小さな方程式系と
有限精度の厳密な非消滅証明によって示す。
単に同じ signature の別成分に例を見つけることでは、この非空性は従わない。

$L_i$ の百二十六係数と対称行列 $A$ の二十八係数を変数とし、
$A_{77}=1$ で規格化する。変数の総数は百五十三である。
次の六十個の方程式を課す。
\begin{align*}
 a_i A B_i^t&=0 &&(i=1,2,3),\\
 a_i A B_j^t+a_j A B_i^t&=0 &&(i<j),\\
 B_i A B_j^t-B_j A B_i^t&=0 &&(i<j).
\end{align*}
最初の二行は三十個、最後の交代行列の独立成分も三十個である。
これらは指定された自己反対 incidence の全方程式であり、
方程式の一部だけを満たす近似を取っているのではない。
基点での pairing は
\[
 A_{12}=A_{21}=1,\quad
 A_{34}=A_{43}=-1,\quad
 A_{56}=A_{65}=1,\quad A_{77}=1
\]
で、他の成分は零である。

付属コードの指定する六十変数に関する Jacobian の行列式は厳密に $2$ である。
残る九十三変数を基点の値に $T\eta$ を加えたものに固定する。
ここで $\eta\in\{-2,-1,0,1,2\}^{93}$ は証明書に明記した整数ベクトルである。
形式的陰関数定理は
\[
 L(T),A(T)\in\mathbb Z_{(101)}[[T]]
\]
に一意な解を与える。
この形式曲線は基点を通る滑らかな自己反対成分に含まれる。
特に、その成分の定義イデアルの各多項式は形式曲線上で恒等的に消える。
基点では左列の独立性、$A$ の可逆性、復元階数四十八、
$\operatorname{rank}F=39$ と $\operatorname{rank}N=36$ の単位小行列を確認している。
従って形式曲線は、前節の $Z_2$ と同じ核枠の開chartに入る。
百五十三変数から六十方程式を除いた次元九十三は、
左関係の基底変換 $\GL_7$ の次元四十九を除くと $\dim Z_2=44$ と一致する。
$T=101$ で収束評価して得る解にも同じ恒等式が残るので、
この解は既存の $Z_2$ と同じ成分上にある。
六十個の残る変数は、多項式系のJacobianが可逆な孤立解を与えるから、この解の座標は
$\mathbb Q$ 上代数的である。
従って $\overline{\mathbb Q}$ の複素埋込みにより複素点も得られる。

\begin{center}
\begin{tabular}{lr}
検証する量&厳密な結果\\\hline
基点の六十次 Jacobian 小行列式&$2$\\
使用素数と計算精度&$p=101$, $p^{1024}$\\
七個の sextic の最高次係数行列 $P$&$v_p(\det P)=9$\\
$P^{-1}$ の分母に必要な $p$ の最大指数&$2$\\
$p^2y$ の乗法行列の特性多項式の判別式&$v_p=700$\\
この乗法行列と判別式の保証精度&$p^{1022}$\\
$q=aAa^t$ の三次対称係数行列の行列式&$v_p=3$
\end{tabular}
\end{center}

ここでの各数値の意味を説明する。
$z=1$ で七個の maximal minors を書き、
次数六の七個の単項式の係数行列を $P$ とする。
$\det P\ne0$ なので無限遠点はなく、すべての六次単項式を
次数五以下の二十一個の単項式に還元できる。
maximal-minor イデアルの高さの上界は二であるため、有限性から高さは二となる。
Hilbert--Burch 分解
\[
 0\longrightarrow\mathcal O(-7)^6\longrightarrow
 \mathcal O(-6)^7\longrightarrow\mathcal I_Z\longrightarrow0
\]
は長さ二十一を与える。従って上の二十一単項式は商代数の基底である。
その基底で $y$ の乗法行列を作り、$p^2$ を掛けて整係数化した。
特性多項式の判別式が非零なので、この商代数は有限エタールである。

精度の損失も有限証明書の一部である。
入力係数の誤差は $p^{1024}$ の倍数であり、
$P^{-1}$ の各成分の valuation は $-2$ 以上である。
逆行列の摂動誤差の valuation は $1024-4$ 以上となり、
$p^2$ を掛けた乗法行列の誤差は $1022$ 以上となる。
特性多項式と判別式は行列成分の整係数多項式なので、同じ精度が保たれる。
計算された valuation $700$ は $1022$ より小さいので、
これは浮動小数点による判定でも、不確かな切断級数の判定でもない。
また $v_p(\det q)=3$ は二次曲線の非退化性を保証する。

以上で命題の仮定を同じ $Z_2$ 上で満たす点が得られた。
これらの条件は開条件なので、その非空開部分で定理の配置が成立する。
この開部分では相対Hilbert--Burch分解と一定のHilbert多項式により
長さ二十一の有限平坦族が得られ、被約条件を加えると有限エタールとなる。
基点 $B_-$ 自身ではmaximal minorsが三角形の方程式 $x_1x_2x_3$ を共通因子に持つ。
従ってここでは、点族の有限平坦性を $B_-$ まで延長しているのではない。
同じ既約成分に属することは関係族の形式曲線によって保証した。
指定対合を持つ有限エタール被覆の monodromy は対合と可換であり、
固定十一点と非固定五対を保つため、主張した群に含まれる。
同型類の次元十一は前節で計算した $44-33$ から従う。

\paragraph{検証ファイルと主張の境界.}
\nolinkurl{verify_self_opposite_geometry_v24.py} が全検証を再実行する。
\nolinkurl{hensel_chart.py} は全六十方程式、基点、Jacobian、
Hensel 反復を記述し、\nolinkurl{hensel_points.py} は長さ、被約性、精度を検証する。
記録された巨大整数を信用する必要はなく、素数、自由方向、精度から再計算できる。
この結果は、$Z_2$ 上の点 monodromy が上記の最大許容群と等しいこと、
指定反対合以外の有限自己同型がないこと、
また「十点が二次曲線上にある」という条件の逆向きの特徴付けまでは主張しない。

% v24 追加用。11+10節の後に置き、共通のHensel/精度証明は参照して再利用する。
\subsection{三種類の署名に対する一つの計数公式}
\label{subsec:all-self-opposite-signatures-v24}

前節の分解は、反対合の一方の固有ブロックが平面上で階数を落とさない場合に
統一して記述できる。

\begin{proposition}[固定点数の署名公式]
\label{prop:self-opposite-signature-count-v24}
$p+q=6$、$1\leq p\leq3$ とし、右平面との同一視後に
\[
 J=\operatorname{diag}(I_p,-I_q),\qquad
 L=\begin{pmatrix}P\\B\end{pmatrix},\qquad R_i=J L_i
\]
と書く。$A$ は対称可逆で、\emph{全ての}自己反対 incidence 方程式
\begin{equation}\label{eq:full-self-opposite-incidence-v24}
 L_i A L_j^tJ=J L_j A L_i^t\qquad(1\leq i,j\leq3)
\end{equation}
を満たすと仮定する。
さらに、$P(x)$ は $\mathbb P^2$ の全点で階数 $p$ を持ち、
$L$ の maximal-minor スキーム $Z$ は有限被約であるとする。
このとき指定反対合による点配置上の対合は
\[
 f(p)=21-12p+2p^2
\]
個の点を固定し、残る $2p(6-p)$ 個を $p(6-p)$ 対に交換する。
非固定点はちょうど
\[
 \Gamma_P=V\bigl(\det(P(x)AP(x)^t)\bigr)
\]
上にあり、固定点はその外にある。
$\det(PAP^t)$ は零多項式ではなく、$\Gamma_P$ は次数 $2p$ の行列式曲線である。
この主張だけから曲線の被約性、既約性、滑らかさは帰結しない。
\end{proposition}
\begin{proof}
全 incidence 方程式により thin module の標準対応が定まり、
前節と同様に、それは $Ju$ を左核として持つ点を対応させる対合である。
$P$ が常に階数 $p$ なので、左核が正固有空間に含まれることはない。
従って固定点は
\[
 T=\{x:\operatorname{rank}B(x)\leq q-1\}
\]
と一致する。
式 \eqref{eq:full-self-opposite-incidence-v24} の混合ブロックから
$P(x)AB(x)^t=0$ を得るので、$B^t$ は
\[
 \mathcal O^q\longrightarrow
 E=\ker\bigl(\mathcal O^7\xrightarrow{\,PA\,}\mathcal O(1)^p\bigr)(1)
\]
を定める。$E$ の階数は $q+1$ である。
$I_6(L)\subset I_q(B)$ より、$T$ は $Z$ の閉部分スキームなので有限被約であり、
局所 Hilbert--Burch 表示、または退化 locus の Chern 類公式により
$\operatorname{length}T=c_2(E)$ となる\cite[\S14.4]{FultonIntersection}。
ねじった完全列
\[
 0\longrightarrow E\longrightarrow\mathcal O(1)^7
  \longrightarrow\mathcal O(2)^p\longrightarrow0
\]
から
\[
 c(E)=\frac{(1+H)^7}{(1+2H)^p},\qquad
 c_2(E)=\bigl(21-14p+2p(p+1)\bigr)H^2=f(p)H^2
\]
を得る。

非固定点では $P,B$ がともに最大階数を持つが、
それらを合わせた $L$ の階数は五である。
従って二つの行空間の共通部分は非零である。
混合ブロックの直交性から、この共通部分は $P$ の Gram 形式の根基に入る。
よって $\det(PAP^t)=0$ となる。

固定点で Gram 行列が特異だと仮定し、その核の非零ベクトルを $c$ とする。
$P$ は全階数なので $k=AP^tc\ne0$ であり、$k\in\ker L$ である。
負固有空間にある左核 $u$ に対して $PAB^t=0$ を微分すると
$u^t(dL)k=0$ を得る。
前節の法線微分と同じ議論で、これは $Z$ の被約性に反する。
従って固定点は $\Gamma_P$ の外にある。
$f(p)>0$ なので、そのような固定点で Gram 行列式は零でなく、
行列式自体も零多項式ではない。
\end{proof}

\begin{remark}[行列式曲線の基底不変性]
復元pairingがスカラー倍を除いて一意な開部分では、
$\Gamma_P$ は指定反対合と全階数の固有ブロックに付随する。
関係枠の変更 $L\mapsto LM$ に対し $A\mapsto M^{-1}AM^{-t}$ なので、
$PAP^t$ は変わらない。
固有ブロックの基底変更 $P\mapsto SP$ では $PAP^t\mapsto S(PAP^t)S^t$ となり、
行列式は非零定数倍されるだけである。
左右平面の座標変更に対しては同じ射影変換で曲線が移る。
従ってこれは任意に選んだ関係枠だけに依存する曲線ではない。
\end{remark}

\begin{theorem}[三つの自己反対族の被約点配置]
\label{thm:all-self-opposite-configurations-v24}
既存の固定成分 $Z_1,Z_2,Z_3$ は、それぞれ次の性質を持つ非空開部分を含む。
各配置は相異なる二十一点からなり、指定反対合による分解は以下のとおりである。
\[
\begin{array}{c|c|c|c|c|c}
 \text{固定成分}&(p,q)&\text{固定点}&\text{交換される対}&
 \text{非固定点を含む曲線の次数}&\text{同型類の次元}\\\hline
 Z_1&(3,3)&3&9&6&7\\
 Z_2&(1,5)&11&5&2&11\\
 Z_3&(2,4)&5&8&4&8
\end{array}
\]
各行で固定点は表示された行列式曲線の外にある。
$Z_2$ の二次曲線は非退化に取れる。
点配置の monodromy は、各行の固定点数を $f$、対の数を $m$ として
\[
 \mathfrak S_f\times\bigl((\mathbb Z/2\mathbb Z)^m\rtimes\mathfrak S_m\bigr)
\]
に含まれる。
\end{theorem}
\begin{proof}
$Z_2$ は前節で証明した。残る二つでも同じ方程式と Hensel 証明を使う。
まず小さい固有ブロックの全階数条件を基点で確認する。

$Z_1$ では正固有ブロックを $u_{12},u_{13},u_{23}$ の三行に取る。
これらの行の一次従属を $(c_{12},c_{13},c_{23})$ とすると
\[
 xc_{12}=yc_{12}=xc_{13}=zc_{13}=yc_{23}=zc_{23}=0,
 \qquad zc_{12}+yc_{13}+xc_{23}=0
\]
を得る。どの $[x:y:z]\in\mathbb P^2$ でも全ての $c_{ij}$ は零となる。

$Z_3$ では反対合の中央作用の全体符号を変える。
これは関係部分空間上の対合を変えない。
基底を
\[
 u_{11}-u_{22},\ u_{13}-u_{23},\
 u_{11}+u_{22},\ u_{13}+u_{23},\ u_{12},\ u_{33}
\]
に変えると、小さい正固有ブロックは最初の二行である。
基点でそれらは
\[
 (0,-x,0,y,-z,z,0),\qquad
 (z,0,z,0,-y,-x,y-x)
\]
となる。
$z\ne0$ なら第零・第二列により独立であり、
$z=0$ なら二つの非零行は互いに交わらない列に支持を持つので独立である。
いずれも全階数条件は開条件である。
有限証明書ではさらに、それぞれの小行列式が
全ての三次単項式、全ての二次単項式を生成することを標数 $101$ で直接確認した。

一般の $(p,q)$ で課す方程式は、混合ブロックの
\[
 P_iAB_i^t=0,\qquad P_iAB_j^t+P_jAB_i^t=0\quad(i<j)
\]
と、\emph{両方の}同符号ブロックの
\[
 P_iAP_j^t-P_jAP_i^t=0,\qquad
 B_iAB_j^t-B_jAB_i^t=0\quad(i<j)
\]
である。独立成分の数は
\[
 6pq+3\left(\binom p2+\binom q2\right)=45+3pq
\]
となる。$p>1$ のとき、正固有ブロックの方程式も省略していない。
変数は前節と同じ百五十三個である。
基点で全方程式の Jacobian は全射であり、その次元は既存の固定成分と一致する。
右外側作用 $G_1=I,G_3=P$ は右座標に吸収し、
元の固定成分との同一視を保つ。

同じ整数方向選択規則と $T=101$ の収束評価を用いた証明書は次の値を与える。
表の $N$ は入力係数の保証精度 $101^N$、$w$ は逆行列の分母指数である。
\[
\begin{array}{c|r|r|r|r|r|r|r}
 &\text{方程式数}&\det(\text{Jacobian})&N&v_{101}(\det P_6)&w&
 v_{101}(\operatorname{Disc})&N-w\\\hline
 Z_1&72&-1&1536&11&3&1120&1533\\
 Z_2&60&2&1024&9&2&700&1022\\
 Z_3&69&12288&1024&9&2&686&1022
\end{array}
\]
ここで $P_6$ は七個の sextic の最高次係数行列であり、
固有ブロック $P$ とは区別する。
判別式は $101^wy$ の乗法行列の特性多項式のものである。
各行でその valuation は保証精度 $N-w$ より小さい。
従って前節と同じ Hilbert--Burch・判別式の議論により、
長さ二十一の有限被約スキームを得る。

基点における $F$、三次関係写像、復元写像の階数は、
いずれの基底でも標数 $101$ でそれぞれ三十九、三十六、四十八である。
従って必要な incidence の開条件も評価後に保たれる。
基点を通る $\mathbb Z_{(101)}[[T]]$ の形式曲線を評価するので、
各例がそれぞれの既存固定成分に属すことも前節と同じく保証される。
命題を適用すれば全ての個数と曲線の次数を得る。
monodromy の包含は指定対合との可換性から従う。
\end{proof}

\paragraph{追加検証ファイル.}
\nolinkurl{hensel_other_signatures.py} と
\nolinkurl{hensel_signature_3_3_v24.json}、
\nolinkurl{hensel_signature_2_4_v24.json} が追加の二例を記録する。
既定の実行は記録済み証明書との読取り専用照合であり、
再生成による書込みは \nolinkurl{--write-certificate} を指定した場合に限る。
既存の $Z_2$ 用方程式を変更せず、両方の同符号ブロックを含む全行列方程式を検証する。

